\section{Case-construction results}
Table~\ref{tab:intake} gives the external qualification counts as of 23 November 2026. The distinction between a scope recommendation and admission matters at this stage. The documentary reader cleared 38 of 160 reviewed findings, deferred 79 for a decision and 25 for additional evidence, and excluded 18. All 38 cleared cases have since reached a build attempt on both historical trees. An earlier snapshot had only 23 build attempts and 11 positive ordinary-use records.

\begin{table}[t]
\caption{External case qualification at the reporting cut. A case is counted once at each stage it reaches. Counts refer to findings, except where explicitly labelled. No row counts agent repairs.}
\label{tab:intake}
\small
\begin{tabular}{@{}p{.57\linewidth}rp{.27\linewidth}@{}}
\toprule
Stage reached & Count & Reference population \\
\midrule
Documentary follow-ups & 333 & 647 finding records \\
Completed benchmark scope reviews & 160 & 333 follow-ups \\
Cleared for technical qualification & 38 & 160 reviews \\
Both historical revisions built & 37 & 38 attempted \\
Specified ordinary use works offline on both & 18 & 38 attempted \\
Findings with a reviewed assessment for A & 9 & 18 offline cases \\
Assessment for A interpretable on both revisions & 8 & 9 reviewed \\
Versioned pages ready for human admission & 5 & 8 with validated A \\
Signed primary admissions & 0 & 5 prepared pages \\
\bottomrule
\end{tabular}
\end{table}

The main loss after compilation occurs at ordinary use. Seventy-four of 76 historical trees built. Only 18 of the 38 findings have positive evidence for affected ordinary use on both revisions offline. Some otherwise runnable tests need historical chain state or another external service. Access to that state would permit further assessment; it would not by itself establish ordinary-use success. Of the nine A assessments reviewed, eight have interpretable results on both revisions and one remains unresolved on the fixed revision. The eight span six upstream families. Five pages across four families are ready for the human reader. A technical pass has not been counted as a human decision.

Evidence for alternative examples is scarcer. The initial scope review cleared no paired cases among the 160. A separate intensive reading revisited twelve candidates across twelve families and judged that two had documentary support for the required distinctions. These are overlapping investigations with different depths of review. Neither of those two cases has signed paired eligibility. In the nine completed A assessments, the existing hidden-evidence review supports one distinct H. A more recent documentary comparison nominated one further candidate H across two instances. Nomination precedes assessment on both references and human sign-off.

\section{Prospective repair experiment}
\label{sec:experiment}
\paragraph{Units and allocation.}
The funded benchmark targets 80 admitted findings in at least 20 upstream families. The allocation includes 40 paired findings in at least 12 families. The configurations are Claude Code with Opus 5.5 and Codex CLI with GPT-5.4-mini, each at medium effort. They are complete repair systems with different tools as well as different models. The exact provider versions, repository images, resource limits and assignment seed are to be frozen before the first session.

Each configuration receives two independent sessions showing the standard example A for each admitted finding. On a paired finding the same configuration receives two additional sessions showing B. Thus the full design has \(80\times2\times2=320\) standard sessions and \(40\times2\times2=160\) additional sessions. The A sessions on a paired finding also contribute to the standard benchmark and are run only once. With twelve paired findings the corresponding count would be 368 and the randomized comparison would be exploratory. The smaller alternative requires a separate decision before seeing repair outcomes.

\paragraph{What changes between assignments.}
A paired case needs two independently justified behaviours from the same defect and a distinct common observation H that both groups will be graded on. The approved examples are designated A and B without consulting repair outcomes. Within each finding and configuration, two of four slots are randomly assigned A and two B. The remainder of the report, repository, ordinary-use suite, allowed tools, effort and time budget are identical within that block. A visible example cannot also occur in the common environment under a second name. Each session starts from the same pinned tree in a fresh workspace. The custodian retains the complete assignment and interleaves run order.

A scorer who cannot see the assignment grades both behaviours on both groups and also records H and ordinary use. A behavioural success requires a build, preserved ordinary use, and a reason-supported correction on that behaviour. The score permits a legitimate implementation different from the reference patch. A session ending in refusal, timeout or an empty patch remains in its assigned group. Infrastructure failures, unresolved assessments and retries retain their own records. First attempts stay visible throughout the analysis.

\paragraph{RQ1 and RQ2.}
For each configuration, let \(W\) indicate a working repair of the reported behaviour A. Average \(W\) over the two assigned A sessions for each finding, then give findings equal weight. Report the raw numerator and denominator and the number lost at patch production, build, ordinary use and A. Also weight families equally and show paired differences between configurations on the same findings. An unresolved assessment earns no demonstrated success in the primary rate; bounds also count unresolved cases at each extreme.

Let \(\mathcal{H}\) contain the primary cases with human-signed, distinct hidden evidence at the roster freeze. Among the same standard A attempts on \(\mathcal{H}\), report both \(W\) and \(W\) together with correction of H. Their difference measures the successes lost when that second obligation is added. Also show how many primary successes remain unevaluable on H, and how the cases in \(\mathcal{H}\) differ in family and reference route. Where H is missing inside the subset, its grade remains unresolved. Outside the subset no H grade is defined. A small difference on a narrow subset says little about other uses of the remaining contracts.

\paragraph{RQ3.}
For a paired finding and configuration, let \(Y_A\) and \(Y_B\) indicate working repairs of the respective behaviours on the same patch. For \(x,z\in\{A,B\}\), write \(\bar Y_{x\mid z}\) for the mean on behaviour \(x\) when example \(z\) was shown. The block contrast is
\begin{equation}
D = \tfrac{1}{2}\left[
  \big(\bar Y_{A\mid A}-\bar Y_{A\mid B}\big)
  +\big(\bar Y_{B\mid B}-\bar Y_{B\mid A}\big)\right].
\label{eq:contrast}
\end{equation}
For example, if every patch corrects A and none corrects B under either report, A is easier on these attempts but both differences are zero. If showing A helps A, while showing B makes little difference to B, the two terms reveal the asymmetry even when their mean is positive. Report both terms as well as their mean, averaged equally across finding-by-configuration blocks.

For the 40-pair design, the primary inference uses the random assignment of example labels. Permute the balanced labels within each block under the sharp no-effect null. Report an assignment-based 95\% interval for the average contrast and identify the large-sample approximation. Also report a bounded-score interval under independent block assignment; it governs the decision if the estimated standard error vanishes. With 80 independent finding-by-configuration blocks its conservative 95\% half-width is about 30 percentage points. With 24 blocks in the exploratory alternative it is about 55 points. These bounds describe the design and are not observed precision. Claim directional support only with a positive estimate, an interval excluding zero, a two-sided randomization \(p<.05\), and a positive lower bound after accounting for unresolved grades. Report the configurations separately, leave one family out at a time, and distinguish external findings from archive findings used in development. A wide interval leaves the question open. If nearly all patches satisfy both behaviours, report the remaining room for improvement under each assignment.

\paragraph{Results outstanding.}
There are no released sessions in the expanded benchmark, so RQ1 has no repair-rate denominator yet. RQ2 has no signed hidden subset, and RQ3 has no signed paired cases or observed assignment contrast. A future results section must show both the actual case flow and all assigned attempts, including unsuccessful and unresolved attempts. No performance estimate or significance test can be filled in from the qualification counts in Table~\ref{tab:intake}.
